\begin{abstract}
Computer-use agents act by choosing among the elements on a screen. The accessibility (AX) trees they usually
rely on omit icons, canvas widgets and custom controls, and pixel parsers that caption every detected element are
slow and hand the agent hundreds of options. We study how to build, from pixels alone, the list of on-screen options
that a \emph{separate} decision model chooses from. Our system proposes options and never clicks, so we ask whether
the needed element is on a short list and how quickly the list is built, not whether one predicted point is
correct. Candidates come from OCR merged with a tiled detector trained on accessibility labels and a public GUI
detector; together they contain the target element for 83.8\%, 96.0\% and 98.5\% of tasks on \osw, \sspro and
\ssv. Labelling candidates costs more than detecting them, so we rank before labelling: a SigLIP model fine-tuned
on context crops with a per-screen listwise loss scores every candidate against the instruction. The top 50
candidates contain the target in 75.1\%, 81.8\% and 98.5\% of tasks (39.5\%, 31.0\% and 93.9\% when ordered by
detector confidence) and need up to 8$\times$ fewer icon captions. Ranked lists complement single-point GUI
grounders: adding a grounder's click to our scores gives lists that contain the target more often than either
alone (\sspro: 89.0\% in the top 50). An off-the-shelf VLM selector given our top 10 instead of the full candidate
list picks the right element 1.5$\times$ more often on \osw and 3.7$\times$ more often on \sspro, from prompts
4--8$\times$ shorter.
\end{abstract}
